\manualchapter{Patches and troubleshooting}{sec:operation}
\subsection{Two colors from one source}
Patch audio and an envelope to the left lane and leave the corresponding
right inputs empty. Set different frequency knobs and route the two outputs
to separate mixer channels. Lower the right Output Level below zero to invert
it before summing externally; compare the resulting cancellation and color.
For conventional stereo processing, use explicit left and right audio cables
and matching settings.

The frequency control changes the interpolation phase increment, with a
response affected by the mode and host sample rate. Tune it by ear; its Hz
label is a pitch conversion and should not be read as a measured cutoff.
Coarse or gritty audio can be the 8-bit conversion rather than unwanted
feedback. Reduce Input Gain to avoid additional clipping.

\subsection{Saved modes and polyphony}
Parameters and the chosen filter mode are saved in patches. Reset selects
Loud, and randomize chooses a mode. The current button updates a loudness
compensation value that the context menu and patch recall do not update.
If the recalled mode has a different level, cycle through the modes with the
button to return to the desired one; this known inconsistency is recorded
for a runtime fix. Filter history is not serialized.

The widest input selects 1--16 channels for both lanes. Inputs read matching
polyphonic channels rather than repeating a mono envelope across all of them.
Match audio and envelope channel counts. The audio path evaluates its controls
every sample; only meter updates are divided down.

For silence, check Input Gain, signed Output Level and its CV. If level or
frequency CV has no effect, check Bypass. When comparing modes, use the same
selection method and allow for level changes before judging the timbre.
